\section{Uczenie agentów jako problem optymalizacji} \label{sec:agent_optim}

Projektowanie skutecznych agentów grających w cyfrowe gry karciane wiąże się z koniecznością doboru wielu parametrów, które wpływają na zachowanie agenta w trakcie rozgrywki. W przypadku agentów opartych na heurystycznych funkcjach oceny stanu planszy, które są powszechnie stosowane w środowisku SabberStone \cite{hearthstonewiki}, parametry te przyjmują postać wag. Wagi te przypisywane są poszczególnym cechom stanu gry, takim jak punkty zdrowia bohaterów, liczba i statystyki stronników na planszy oraz dostępne zasoby many. Znalezienie zestawu wag, który maksymalizuje skuteczność agenta mierzoną wskaźnikiem wygranych, jest z definicji zadaniem optymalizacji numerycznej \cite{sanchez2020}.

Kolejną trudnością jest to, że problem ten ma charakter czarnej skrzynki (ang. \textit{black-box optimization}). Oznacza to, że nie da się wyrazić relacji między konkretnymi wartościami wag a skutecznością agenta w postaci formuły analitycznej. Ocena ta wymaga przeprowadzenia rzeczywistych symulacji rozgrywek w środowisku gry. Skutkuje to również tym, że gradientowe metody optymalizacji nie mają tutaj zastosowania. Dodatkowym utrudnieniem jest szum funkcji celu, który wynika z losowych efektów gry --- ten sam zestaw parametrów oceniany w różnych rozgrywkach może dawać odmienne wyniki \cite{hoover2020}. Tego rodzaju problemy optymalizacji stochastycznej są szczególnie dobrze zbadane w kontekście algorytmów ewolucyjnych \cite{TogeliusYannakakis2023, TogeliusYannakakis2025}.

Osobnym problemem optymalizacyjnym jest budowanie talii (ang. deck building). Jak wykazali Yang i in. \cite{yang2021}, dobór 30 kart spośród puli liczącej kilkaset pozycji, z ograniczeniami co do liczby kopii i klasy bohatera, tworzy ogromną kombinatoryczną przestrzeń poszukiwań, której nie sposób przeszukać metodami wyczerpującymi. Kowalski i Miernik \cite{kowalski2023} potwierdzają, że optymalizacja talii w warunkach konkursowych stanowi wyzwanie porównywalne pod względem trudności z optymalizacją samej strategii gry. Oba zadania --- optymalizacja wag funkcji oceny oraz budowanie talii --- są zatem naturalnymi kandydatami do rozwiązywania metodami ewolucyjnymi \cite{TogeliusYannakakis2023}.

